\chapter{Loops}

\section{\texttt{for}}

Iterates over a sequence. The iteration variable is visible inside the
block and \textit{persists} after the loop with the last value it held:

\begin{lstlisting}[caption={for over a range}]
for i in 1..4 {
  display i    // 1, 2, 3  (4 not included)
}
display i    // 3 — variable persists after the loop
\end{lstlisting}

\subsection*{Exclusive and inclusive ranges}

\hayashi{} follows the same convention as Rust:

\begin{lstlisting}
for i in 1..3  { }    // 1, 2      (exclusive: end not included)
for i in 1..=3 { }    // 1, 2, 3   (inclusive: end included)
\end{lstlisting}

Ranges are also expressions that evaluate to lists — see
Section~\ref{sec:ranges}.

\subsection*{Iterating over lists}

\begin{lstlisting}
let fruits = ["apple", "banana", "grape"]

for fruit in fruits {
  display fruit
}
\end{lstlisting}

\subsection*{Discarding the iteration variable}

When the index is not needed:

\begin{lstlisting}
for _ in 1..3 {
  display "repetition"
}
\end{lstlisting}

\section{\texttt{while}}

Executes while the condition is true:

\begin{lstlisting}[caption={Basic while}]
let n = 1

while n <= 5 {
  display n
  n += 1
}
\end{lstlisting}

\begin{aviso}
  \hayashi{} does not detect infinite loops. A \lstinline|while true { }|
  without \lstinline{break} will run indefinitely. Use \texttt{Ctrl-C}
  to interrupt.
\end{aviso}

\section{\texttt{break} and \texttt{continue}}

\lstinline{break} exits the loop immediately. \lstinline{continue} advances
to the next iteration:

\begin{lstlisting}[caption={break and continue}]
for i in 1..10 {
  if i == 3 { continue }    // skip 3
  if i == 7 { break }       // stop at 7
  display i
}
// prints: 1, 2, 4, 5, 6
\end{lstlisting}

\section{Loops with DataFrames}

\lstinline{for} can iterate over lists produced by data functions:

\begin{lstlisting}[caption={Iterating over groups}]
load "sales.csv" as df

for state in ["TX", "CA", "NY"] {
  let sub = filter(df, region == state)
  display f"{state}: {count(sub)} obs"
}
\end{lstlisting}

\section{\texttt{parallel for}}

Concurrent variant of \texttt{for}. Iterations run in parallel across
threads; each iteration's return value (last expression or explicit
\texttt{return}) is collected into a list, which is stored back into the
\textbf{iteration variable} after all threads complete.

\begin{lstlisting}[caption={Basic parallel for}]
parallel for t in tickers {
  let sub = filter(df, ticker == t)
  nrow(sub)
}
// t now holds a list of nrow values, one per ticker
\end{lstlisting}

Optional \texttt{threads=N} limits the number of worker threads
(default: all available CPUs):

\begin{lstlisting}[caption={Thread limit}]
parallel for t in tickers, threads=4 {
  load "data.db" as sub, query=f"SELECT * FROM p WHERE t='{t}'"
  nrow(sub)
}
\end{lstlisting}

Each thread gets its own interpreter with a snapshot of the outer
environment (only send-safe values are captured). Use \texttt{return nil}
to skip an iteration; \texttt{nil} entries are kept in the result list.

Combine with \texttt{rbind()} to aggregate per-iteration DataFrames:

\begin{lstlisting}[caption={parallel for + rbind}]
parallel for i in 0..n, threads=8 {
  let t = tickers[i]
  if t == "SPY" { return nil }
  let df_t = compute_betas(t)
  df_t
}
let all = rbind(i)   // i holds the list of DataFrames
\end{lstlisting}
